\documentclass[a4paper,12pt]{article}

% ============================================================
% Общий преамбул для отчётов по лабораторным работам (labs_lpr)
% Подключается из report/main.tex командой:
%   % ============================================================
% Общий преамбул для отчётов по лабораторным работам (labs_lpr)
% Подключается из report/main.tex командой:
%   % ============================================================
% Общий преамбул для отчётов по лабораторным работам (labs_lpr)
% Подключается из report/main.tex командой:
%   \input{../../common/latex_template/preamble.tex}
% ============================================================

% Поддержка русского языка и кодировки
\usepackage[T2A]{fontenc}
\usepackage{lmodern}
\usepackage[utf8]{inputenc}
\usepackage[russian]{babel}

% Математические пакеты
\usepackage{amsmath, amssymb}

% Настройка полей
\usepackage{geometry}
\geometry{top=2cm, bottom=2cm, left=2.5cm, right=1.5cm}

% Колонтитулы
\usepackage{fancyhdr}
\usepackage{graphicx}

% Современные таблицы
\usepackage{tabularray}

% Работа с csv файлами
\usepackage{csvsimple-l3}
\usepackage{tikz}

% Графики
\usepackage{pgfplots}
\pgfplotsset{compat=1.18}

% Ссылки
% \usepackage{hyperref} % гипперссылки, нужно оформление вместо красного квадрата
\usepackage[russian]{cleveref} % умные ссылки, нужно переоформление для формул, рисунков и тд (вместо ф-л что-то нужно придумать)

% ============================================================

% Поддержка русского языка и кодировки
\usepackage[T2A]{fontenc}
\usepackage{lmodern}
\usepackage[utf8]{inputenc}
\usepackage[russian]{babel}

% Математические пакеты
\usepackage{amsmath, amssymb}

% Настройка полей
\usepackage{geometry}
\geometry{top=2cm, bottom=2cm, left=2.5cm, right=1.5cm}

% Колонтитулы
\usepackage{fancyhdr}
\usepackage{graphicx}

% Современные таблицы
\usepackage{tabularray}

% Работа с csv файлами
\usepackage{csvsimple-l3}
\usepackage{tikz}

% Графики
\usepackage{pgfplots}
\pgfplotsset{compat=1.18}

% Ссылки
% \usepackage{hyperref} % гипперссылки, нужно оформление вместо красного квадрата
\usepackage[russian]{cleveref} % умные ссылки, нужно переоформление для формул, рисунков и тд (вместо ф-л что-то нужно придумать)

% ============================================================

% Поддержка русского языка и кодировки
\usepackage[T2A]{fontenc}
\usepackage{lmodern}
\usepackage[utf8]{inputenc}
\usepackage[russian]{babel}

% Математические пакеты
\usepackage{amsmath, amssymb}

% Настройка полей
\usepackage{geometry}
\geometry{top=2cm, bottom=2cm, left=2.5cm, right=1.5cm}

% Колонтитулы
\usepackage{fancyhdr}
\usepackage{graphicx}

% Современные таблицы
\usepackage{tabularray}

% Работа с csv файлами
\usepackage{csvsimple-l3}
\usepackage{tikz}

% Графики
\usepackage{pgfplots}
\pgfplotsset{compat=1.18}

% Ссылки
% \usepackage{hyperref} % гипперссылки, нужно оформление вместо красного квадрата
\usepackage[russian]{cleveref} % умные ссылки, нужно переоформление для формул, рисунков и тд (вместо ф-л что-то нужно придумать)


% Настройка колонтитулов
\pagestyle{fancy}
\fancyhf{}
% Верхний левый: логотип (если файла нет, выводится текст)
\fancyhead[L]{%
    \IfFileExists{../../common/latex_template/miptlogo.png}{%
        \includegraphics[height=1cm]{../../common/latex_template/miptlogo.png}%
    }{%
        \textbf{[ЛОГОТИП МФТИ]}%
    }%
}
% Верхний правый: инициалы
\fancyhead[R]{\textbf{Емельянов ДА, Байдин ФМ}}
% Нижний центральный: номер страницы
\fancyfoot[C]{\thepage}
\renewcommand{\headrulewidth}{0.4pt}

\begin{document}

% --- ТИТУЛЬНЫЙ ЛИСТ ---
\begin{titlepage}
\begin{center}
    {\large МОСКОВСКИЙ ФИЗИКО-ТЕХНИЧЕСКИЙ ИНСТИТУТ}
\end{center}
\begin{center}
    {\largeФизтех-школа физики и исследований им. Ландау\break ЛФИ}
\end{center}


    \vspace{3cm}

\begin{center}
   
\end{center}
\vspace{0.15cm}
{
\begin{center}
    \huge \bf Лабораторная работа №3.2.6 по курсу Общая физика\\
    \vspace{0.4cm}      
    \LARGE Изучение гальванометра.
\end{center}
}
\vspace{1.5cm}
\begin{flushright}  
Емельянов Дмитрий Алексеевич\\
Байдин Фаддей Михайлович\\
\vspace{0.4cm}
Б02-502
\end{flushright}
\vspace{10cm}
\begin{center}
    \large 8 сентября 2026, г. Долгопрудный
\end{center}
\end{titlepage}
% ----------------------

\section{Аннотация}
В данной работе проводится изучение работы высокочувствительного зеркального гальванометра магнитоэлектрической системы. Исследуются два режима работы прибора: режим измерения постоянного тока (стационарный) и режим измерения электрического заряда (баллистический). Определяются ключевые характеристики гальванометра: динамическая постоянная, баллистическая постоянная и критическое сопротивление цепи.

\section{Теоретические сведения}
Уравнение движения рамки гальванометра в магнитном поле имеет вид:
\begin{equation}
    \label{eq:motion}
    \ddot{\varphi} + 2\gamma\dot{\varphi} + \omega_0^2\varphi = KI,
\end{equation}
где $\varphi$ --- угол поворота рамки, $\gamma$ --- коэффициент затухания, $\omega_0$ --- собственная частота колебаний, $K$ --- коэффициент, зависящий от параметров гальванометра.

В стационарном режиме при пропускании постоянного тока $I$ угол поворота $\varphi$ связан с током через динамическую постоянную $C_I$:
\begin{equation}
    \label{eq:ci_def}
    C_I = \frac{I}{\varphi} = \frac{2aI}{x},
\end{equation}
где $a$ --- расстояние от шкалы до зеркальца, $x$ --- координата светового пятна на шкале. 

Сила тока в цепи рассчитывается по формуле:
\begin{equation}
    \label{eq:current}
    I = \frac{R_1}{R_2}\frac{U_0}{R+R_0},
\end{equation}
где $U_0$ --- показания вольтметра, $R_1/R_2$ --- положение делителя, $R$ --- сопротивление магазина, $R_0$ --- внутреннее сопротивление гальванометра.

Для режима свободных затухающих колебаний логарифмический декремент затухания $\Theta$ равен:
\begin{equation}
    \label{eq:decrement}
    \Theta = \ln\frac{x_n}{x_{n+1}},
\end{equation}
где $x_n$ и $x_{n+1}$ --- два последовательных отклонения зайчика в одну сторону.

Критическое сопротивление цепи гальванометра $R_{kp}$ связано с $\Theta$ следующим образом:
\begin{equation}
    \label{eq:rkp_formula}
    R_{kp} = \frac{R+R_0}{\sqrt{\left(\frac{2\pi}{\Theta}\right)^2+1}} - R_0.
\end{equation}

В баллистическом режиме через гальванометр пропускается заряд $q = C U_C$. Баллистическая постоянная в критическом режиме равна:
\begin{equation}
    \label{eq:cq}
    C_q^{kp} = \frac{q}{\varphi_{max}^{kp}} = 2a\frac{R_1}{R_2}\frac{C U_0}{x_{max}^{kp}}.
\end{equation}

Для гальванометра, работающего в баллистическом режиме с сопротивлением $R > R_{kp}$ (слабое затухание), первый отброс $l$ связан с отбросом разомкнутого гальванометра $l_0$ соотношением
\begin{equation}
    \label{eq:l0_l}
    l_0 = l\, e^{\Theta/4}
    \quad\Rightarrow\quad
    \Theta = 4\ln\frac{l_0}{l}.
\end{equation}
Поскольку коэффициент затухания обратно пропорционален сопротивлению цепи, $\gamma \propto 1/(R+R_0)$, декремент связан с критическим сопротивлением формулой
\begin{equation}
    \label{eq:lin_theory}
    \frac{4\pi^2}{\Theta^2} + 1 = \frac{(R+R_0)^2}{(R_{kp}+R_0)^2}.
\end{equation}
Введя переменные $X = (R+R_0)^2$ и $Y = 1/\Theta^2$, получаем линейную зависимость
\begin{equation}
    \label{eq:lin}
    Y = kX + b,\qquad k = \frac{1}{4\pi^2 (R_{kp}+R_0)^2}
    \quad\Rightarrow\quad
    R_{kp} = \frac{1}{2\pi\sqrt{k}} - R_0.
\end{equation}

\section{Используемое оборудование}
\begin{itemize}
    \item Зеркальный гальванометр с осветителем и шкалой;
    \item Источник постоянного напряжения;
    \item Делитель напряжения;
    \item Магазин сопротивлений;
    \item Эталонный конденсатор;
    \item Вольтметр;
    \item Переключатель, ключи, линейка.
\end{itemize}

\section{Методика измерений}
Измерения проводятся в двух конфигурациях цепи. Для стационарного режима используется подключение гальванометра к делителю напряжения. Регулировка сопротивления магазина $R$ позволяет изменять ток в цепи и переводить гальванометр из колебательного режима в апериодический. Измеряя отклонения зайчика, можно рассчитать динамическую постоянную по формуле \cref{eq:ci_def}.

В баллистическом режиме цепь содержит эталонный конденсатор $C$, который заряжается от источника и затем разряжается через гальванометр, сообщая ему короткий импульс тока. Первый максимальный отброс $x_{max}$ используется для определения баллистической постоянной $C_q^{kp}$ согласно \cref{eq:cq}. По зависимости отброса от сопротивления с помощью \cref{eq:l0_l,eq:lin} линеаризацией определяется критическое сопротивление.

\section{Результаты измерений и обработка данных}

\subsection{Параметры экспериментальной установки}
\begin{itemize}
    \item Внутреннее сопротивление гальванометра: $R_0 = 560$ Ом;
    \item Напряжение источника: $U_0 = \frac{67}{150}\cdot3\text{ В}\approx1,34$ В;
    \item Расстояние от зеркальца до шкалы: $a = 0,98$ м;
    \item Положение делителя напряжения: $R_1/R_2 = 1 / 2000$;
    \item Ёмкость эталонного конденсатора: $C = 2$ мкФ.
\end{itemize}

\subsection{Определение динамической постоянной (стационарный режим)}

Результаты измерений отклонения зайчика в зависимости от сопротивления магазина представлены в \cref{tab:stationary_data}. Сила тока $I$ рассчитана по формуле \cref{eq:current}.

\begin{table}[!ht]
    \centering
    \caption{Измерения в стационарном режиме}
    \label{tab:stationary_data}

    \csvreader[
        tabularray = {
            colspec = {|c|c|c|c|c|},
            hlines,
            vlines,
            row{1} = {font=\bfseries},
        },
        table head = {
            \textbf{$R$, кОм}
            &
            \textbf{$x$, см}
            &
            \textbf{$\sigma_x$, см}
            &
            \textbf{$I$, мкА}
            &
            \textbf{$\sigma_I$, мкА}
            \\
        }
    ]%
    {../data/raw/data_stationary.csv}%
    {x=\csvX, sigma_x=\csvSigmaX, R=\csvR, I=\csvI, sigma_I=\csvSigmaI}%
    {%
        $\csvR$
        &
        $\csvX$
        &
        $\csvSigmaX$
        &
        $\csvI$
        &
        $\csvSigmaI$
    }

\end{table}

На основе данных из \cref{tab:stationary_data}
построен график зависимости $I(x)$ с учётом погрешностей измерений
(\cref{fig:ix_graph}).

\begin{figure}[!ht]
    \centering
    \begin{tikzpicture}
        \begin{axis}[
            width=0.8\textwidth,
            height=0.5\textwidth,
            xlabel={$x$, cм},
            ylabel={$I$, мкА},
            enlarge x limits=0,
            enlarge y limits=0,
            grid=both,
            legend pos=north west,
        ]

            \addplot[
                only marks,
                mark=*,
                /pgfplots/error bars/x dir=both,
                /pgfplots/error bars/x explicit,
                /pgfplots/error bars/y dir=both,
                /pgfplots/error bars/y explicit,
            ]
            table[
                col sep=comma,
                x=x,
                y=I,
                x error=sigma_x,
                y error=sigma_I,
            ] {../data/raw/data_stationary.csv};

            \addlegendentry{Экспериментальные данные}

            % Линейная аппроксимация
            \addplot[
                no marks,
                domain=0:25,
                thick
            ] {0.0051 * x + 0.0017};

            \addlegendentry{Линейная аппроксимация}
            
        \end{axis}
    \end{tikzpicture}
    \caption{Зависимость силы тока от отклонения зайчика}
    \label{fig:ix_graph}
\end{figure}

По угловому коэффициенту линейной аппроксимации рассчитываем динамическую постоянную:
\begin{equation}
    C_I = 1,005\cdot10^{-9} \text{ А/(мм/м)}.
\end{equation}
Чувствительность гальванометра к току $S_I = 1/C_I = 995\cdot10^6$ (мм/м)/А.

\subsection{Определение критического сопротивления}

Период свободных колебаний рамки $T_0 = 4,64$ с.
Логарифмический декремент затухания разомкнутого гальванометра $\Theta_0 = 0,249$.

Измерения последовательных отклонений и рассчитанный логарифмический декремент затухания $\Theta$ по формуле \cref{eq:decrement} представлены в \cref{tab:decrement_data}. Столбец $R_{kp}$ рассчитан скриптом \texttt{scripts/process\_data.py} из сырых данных \texttt{data/raw/data\_table.csv} и сохранён в \texttt{data/processed/data\_table.csv}.

\begin{table}[ht]
    \centering
    \caption{Определение декремента затухания}
    \label{tab:decrement_data}

    \csvreader[
        tabularray = {
            colspec = {|c|c|c|c|c|},
            hlines,
            vlines,
            row{1} = {font=\bfseries},
        },
        table head = {
            \textbf{$R$, кОм}
            &
            \textbf{$x_n$, см}
            &
            \textbf{$x_{n+1}$, см}
            &
            \textbf{$\Theta$}
            &
            \textbf{$R_{kp}$, кОм}
            \\
        }
    ]%
    {../data/processed/data_table.csv}%
    {R=\csvR, x1=\csvXn, x2=\csvXnp, Theta=\csvTheta, R_cr=\csvRcr}%
    {%
        $\csvR$
        &
        $\csvXn$
        &
        $\csvXnp$
        &
        $\csvTheta$
        &
        $\csvRcr$
    }

\end{table}

Критическое сопротивление, полученное путём подбора (когда зайчик не переходит за ноль): $R_{kp}^{(\textrm{подбор})} = 7,5\pm0,1$ кОм.

Усредненное значение критического сопротивления по результатам расчётов (\cref{eq:rkp_formula}):
\begin{equation}
    R_{kp} = 8.280 \pm 0.648 \text{ кОм}.
\end{equation}

\subsection{Баллистический режим}

Отброс зайчика в режиме свободных колебаний ($R = \infty$): $x_0 = 245$ мм. При $R_1/R_2 = 1/50$. Зафиксируем $R_1/R_2$ для баллистического режима.
Максимальный отброс гальванометра без затухания: 
\begin{equation}
    x_{max}^{cB} = x_0 \left(1 + \frac{\Theta_0}{4}\right) = 260 \text{ мм}.
\end{equation}

Результаты измерения первого отброса $x_{max}$ при различных сопротивлениях $R$ приведены в \cref{tab:ballistic_data}.

\begin{table}[ht]
    \centering
    \caption{Измерения в баллистическом режиме}
    \label{tab:ballistic_data}

    \csvreader[
        tabularray = {
            colspec = {|c|c|c|c|},
            hlines,
            vlines,
            row{1} = {font=\bfseries},
        },
        table head = {
            \textbf{$R$, кОм}
            &
            \textbf{$R+R_0$, кОм}
            &
            \textbf{$x_{max}$, см}
            &
            \textbf{$\sigma_{x_{max}}$, см}
            \\
        }
    ]%
    {../data/raw/ballistic_mode_data.csv}%
    {1=\csvR, 2=\csvXmax, 3=\csvSigmaXmax, 4=\csvRtotal}%
    {%
        $\csvR$
        &
        $\csvRtotal$
        &
        $\csvXmax$
        &
        $\csvSigmaXmax$
    }

\end{table}

\begin{figure}[!ht]
    \centering
    \begin{tikzpicture}
        \begin{axis}[
            width=0.8\textwidth,
            height=0.5\textwidth,
            xlabel={$R+R_0$, кОм},
            ylabel={$x_{max}$, см},
            enlarge x limits=0.05,
            enlarge y limits=0.05,
            grid=both,
            legend pos=south east,
        ]

            \addplot[
                only marks,
                mark=*,
                /pgfplots/error bars/y dir=both,
                /pgfplots/error bars/y explicit,
            ]
            table[
                col sep=comma,
                x index=3,
                y index=1,
                y error index=2,
            ] {../data/raw/ballistic_mode_data.csv};

            \addlegendentry{Экспериментальные данные}

            \addplot[no marks, domain=3:15, thick] {9.565};
            
            \addlegendentry{$x_{max}^{kp} = x_{max}^{cB}/e \approx 9.6$ см}

        \end{axis}
    \end{tikzpicture}
    \caption{Зависимость первого отброса от суммарного сопротивления}
    \label{fig:ballistic_graph}
\end{figure}

По графику (\cref{fig:ballistic_graph}) определяем $R_{kp}$, при котором $x_{max}$ в $e$ раз меньше $x_{max}^{cB}$:
\begin{equation}
    R_{kp}^{(\text{балл})} = 6,2 \text{ кОм}.
\end{equation}

\subsection{Определение критического сопротивления линеаризацией баллистической кривой}

Для каждого значения $R$ по \cref{eq:l0_l} рассчитан логарифмический декремент затухания, где $l_0 = 24,5\pm0,5$ см --- отброс разомкнутого гальванометра. Погрешности оценивались по формулам
\begin{equation}
    \sigma_\Theta = 4\sqrt{\left(\frac{\sigma_l}{l}\right)^2 + \left(\frac{\sigma_{l_0}}{l_0}\right)^2},
    \qquad
    \sigma_Y = \frac{2\sigma_\Theta}{\Theta^3}.
\end{equation}
Результаты расчёта, выполненного скриптом \texttt{scripts/process\_data.py}, приведены в \cref{tab:linearization}.

\begin{table}[!ht]
    \centering
    \caption{Линеаризация баллистической кривой}
    \label{tab:linearization}
    \small
    \csvreader[
        tabularray = {
            colspec = {|c|c|c|c|c|c|c|},
            hlines,
            vlines,
            row{1} = {font=\bfseries},
        },
        table head = {
            \textbf{$R$, кОм}
            &
            \textbf{$R+R_0$, кОм}
            &
            \textbf{$l$, см}
            &
            \textbf{$\Theta$}
            &
            \textbf{$\sigma_\Theta$}
            &
            \textbf{$X$, кОм$^2$}
            &
            \textbf{$Y \pm \sigma_Y$}
            \\
        }
    ]%
    {../data/processed/ballistic_linearization.csv}%
    {R=\csvR, RR0=\csvRRo, l=\csvL, Theta=\csvTheta, sigma_Theta=\csvSTheta, X=\csvX, Y=\csvY, sigma_Y=\csvSY}%
    {%
        $\csvR$
        &
        $\csvRRo$
        &
        $\csvL$
        &
        $\csvTheta$
        &
        $\csvSTheta$
        &
        $\csvX$
        &
        $\csvY \pm \csvSY$
    }
\end{table}

По данным \cref{tab:linearization} построен график $Y(X)$ (\cref{fig:linearization}). Методом наименьших квадратов получены параметры прямой $Y = kX + b$:
\begin{equation}
    k = (8,96 \pm 0,03)\cdot10^{-4}\ \text{кОм}^{-2},
    \qquad
    b = 0,039 \pm 0,003.
\end{equation}

\begin{figure}[!ht]
    \centering
    \begin{tikzpicture}
        \begin{axis}[
            width=0.8\textwidth,
            height=0.55\textwidth,
            xlabel={$X = (R+R_0)^2$, кОм$^2$},
            ylabel={$Y = 1/\Theta^2$},
            xmin=0, xmax=2700,
            ymin=0,
            grid=both,
            legend pos=north west,
            scaled x ticks=false,
            xtick distance=500,
        ]

            \addplot[
                only marks,
                mark=*,
                mark size=1.5pt,
                /pgfplots/error bars/y dir=both,
                /pgfplots/error bars/y explicit,
            ]
            table[
                col sep=comma,
                x=X,
                y=Y,
                y error=sigma_Y,
            ] {../data/processed/ballistic_linearization.csv};

            \addlegendentry{Экспериментальные данные}

            \addplot[
                no marks,
                domain=0:2700,
                samples=2,
                thick
            ] {8.9610e-4 * x + 0.03945};

            \addlegendentry{МНК: $Y = 8{,}96\cdot10^{-4}X + 0{,}039$}

        \end{axis}
    \end{tikzpicture}
    \caption{Линеаризованная зависимость $1/\Theta^2$ от $(R+R_0)^2$}
    \label{fig:linearization}
\end{figure}

По угловому коэффициенту с помощью \cref{eq:lin} находим критическое сопротивление:
\begin{equation}
    R_{kp}^{(\text{лин})} = \frac{1}{2\pi\sqrt{k}} - R_0 = 4,76 \pm 0,01 \text{ кОм},
\end{equation}
где указана только случайная погрешность, определяемая разбросом точек относительно прямой, $\sigma_{R_{kp}} = \sigma_k / (4\pi k^{3/2})$.

Полученное значение заметно меньше найденного подбором ($7,5$ кОм). Это объясняется тем, что при визуальном подборе глаз фиксирует прекращение колебаний раньше реального момента, а также тем, что соотношение \cref{eq:l0_l} справедливо лишь при слабом затухании ($\Theta \ll 2\pi$), тогда как точки вблизи критического режима имеют $\Theta \approx 3$--$5$. Кроме того, свободный член $b = 0,039$ отличается от теоретического $b = -1/(4\pi^2) \approx -0,025$, что указывает на систематическое отклонение модели.

Расчёт баллистической постоянной в критическом режиме (\cref{eq:cq}):
\begin{equation}
    C_q^{kp} = 2a\frac{R_1}{R_2}\frac{C U_0}{x_{max}^{kp}} = \frac{5,36\cdot10^{-8}\text{ Кл}}{2\cdot 95,7\text{ мм}/(2\cdot 0,98\text{ м})} = 1,097\cdot10^{-9} \text{ Кл/(мм/м)}.
\end{equation}

Время релаксации цепи $\tau = R_0 C = 1,12$ мс.
В сравнении с $T_0 = 4,64$ с: $\tau \ll T_0$, что подтверждает применимость баллистического режима.

\section{Обзор результатов}
Критическое сопротивление было определено четырьмя способами:
\begin{enumerate}
    \item подбором: $R_{kp} = 7,5 \pm 0,1$ кОм;
    \item по логарифмическому декременту затухания: $R_{kp} = 8,3 \pm 0,6$ кОм;
    \item по баллистической кривой (отброс в $e$ раз меньше свободного): $R_{kp} = 6,2$ кОм;
    \item линеаризацией баллистической кривой: $R_{kp} = 4,76 \pm 0,01$ кОм.
\end{enumerate}
Значения, полученные первыми тремя способами, согласуются между собой в пределах $\sim 25\%$. Значение, полученное линеаризацией, занижено: использованное приближение слабого затухания \cref{eq:l0_l} нарушается вблизи критического режима, а указанная погрешность учитывает лишь статистический разброс точек и не отражает систематику модели.

Результаты эксперимента приблизительно сходятся с ожидаемыми значениями. Зависимость $I(x)$ в исследованном диапазоне близка к линейной, что соответствует предположениям модели. Основной вклад в погрешность вносят погрешности измерения $x$ и $I$.

\section{Заключение}
В ходе работы был исследован высокочувствительный зеркальный гальванометр. 
Определены его основные характеристики:
\begin{itemize}
    \item Динамическая постоянная: $C_I = 1,005\cdot10^{-9}$ А/(мм/м);
    \item Баллистическая постоянная: $C_q^{kp} = 1,097\cdot10^{-9}$ Кл/(мм/м);
    \item Критическое сопротивление: $R_{kp} \approx 6$--$8$ кОм (подбор, декремент затухания, баллистическая кривая); линеаризация баллистической кривой даёт $R_{kp} = 4,76 \pm 0,01$ кОм.
\end{itemize}
Показано, что шкала гальванометра обладает высокой степенью линейности, а условие баллистического режима ($\tau \ll T_0$) строго выполняется.

\end{document}
